\scanpage{11}
\mspage{8}
où $\delta$ est le discriminant du second membre de (3.1). Le discriminant $\delta$ étant invariant par translation, il suffit de vérifier cette identité algébrique pour 3 inversible et $a_2=0$.

Lorsque $c_4$ est inversible, un et un seul choix de $x$ et $y$ donne $d_6=0$ \unsure{?|,} i.e. $a_2a_4-9a_6=0$. On a, lorsque cette condition est remplie,
\[
-c_6=2^6(a_2^3-3^3a_6)\qquad \Delta=2^6(-a_2^3a_6-a_4^3+2\cdot3^3\,a_6^2)\tag*{(3.3)}
\]
et
\[
\begin{aligned}
a_2c_4&=-2^{-2}c_6\\
a_4c_4^2&=-2^2\,3^2\,\Delta\\
a_6c_4^3&=c_6\,\Delta\qquad .
\end{aligned}\tag*{(3.4)}
\]

Il suffit de vérifier ces identités en caractéristique 0 , auquel cas la troisième est le produit des deux premières, puisque $d_6=0$.

\uline{Proposition} 3.5.\qquad \uline{Au dessus de} $\mathrm{Spec}(\mathbb Z[2^{-1}])$, \uline{il existe un schéma de module pour les courbes de genre un munies d'une forme différentielle invariante inversible et telles que} $c_4$ \uline{soit inversible. Ce schéma est}
\[
\mathrm{Spec}(\mathbb Z[2^{-1}][a_2,a_4,a_6,(a_2^2-3a_4)^{-1}]/(a_2a_4-9a_6))
\]
\uline{avec pour courbe universelle} (\uline{en coordonnées non homogène}) \handwritten{\tikz[baseline=-.2ex,x=1ex,y=1ex]{\draw (0,0)--(.26,1.20) .. controls (.13,1.47) and (.55,1.58) .. (.47,1.34) .. controls (.42,1.20) and (.29,1.16) .. (.26,1.20);}}\inserted{\handwritten{section $e$}}
\[
y^2=x^3+a_2x^2+a_4x+a_6
\]
% archive sheet number: 10
